\chapter{2015年上海卷（秋理）}

\section{填空题}

\begin{problem}
设全集 \(U = \R\). 若集合 \(A = \{1, 2, 3, 4\}\)，\(B = \{x \mid 2 \leqslant x \leqslant 3\}\)，则 \(A \cap \complement_U B = \)\fillinblank{}.
\end{problem}

\begin{problem}
若复数 \(z\) 满足 \(3z + \overline{z} = 1 + \i\)，其中 \(\i\) 为虚数单位，则 \(z =\fillinblank{}\).
\end{problem}

\begin{problem}
若线性方程组的增广矩阵为 \(\left( \begin{array}{ccc}2 & 3 & c_1\\ 0 & 1 & c_2 \end{array} \right)\) ，解为 \(\begin{cases}x = 3,\\ y = 5 \end{cases}\)， 则 \(c_{1} - c_{2} =\fillinblank{}\)
\end{problem}

\begin{problem}
若正三棱柱的所有棱长均为 \(a\)，且其体积为 \(16\sqrt{3}\)，则 \(a = \)\fillinblank{}.
\end{problem}

\begin{problem}
抛物线 \( y^{2} = 2px  \) （\(p > 0\)）上的动点 \(Q\) 到焦点的距离的最小值为 1，则 p =\fillinblank{}.
\end{problem}

\begin{problem}
若圆锥的侧面积与过轴的截面面积之比为 \(2\pi\)，则其母线与轴的夹角的大小为\fillinblank{}.
\end{problem}

\begin{problem}
方程 \(\log_2(9^{x - 1} - 5) = \log_2(3^{x - 1} - 2) + 2\) 的解为\fillinblank{}.
\end{problem}

\begin{problem}
在报名的 3 名男教师和 6 名女教师中，选取 5 人参加义务献血. 要求男，女教师都有，则不同的选取方式的种数为\fillinblank{}. （结果用数值表示）
\end{problem}

\begin{problem}
已知点 \(P\) 和 \(Q\) 的横坐标相同，\(P\) 的纵坐标是 \(Q\) 的纵坐标的 2 倍，\(P\) 和 \(Q\) 的轨迹分别为双曲线 \(C_1\) 和 \(C_2\). 若 \(C_1\) 的渐近线方程为 \(y = \pm \sqrt{3} x\)，则 \(C_2\) 的渐近线方程为\fillinblank{}.
\end{problem}

\begin{problem}
设 \( f^{-1}(x) \) 为 \(f(x) = 2^{x - 2} + \frac{x}{2}\)，\(x \in [0,2]\) 的反函数，则 \( y = f(x) + f^{-1}(x) \) 的最大值为\fillinblank{}.
\end{problem}

\begin{problem}
在 \(\left(1 + x + \frac{1}{x^{2015}}\right)^{10}\) 的展开式中，\(x^2\) 项的系数为\fillinblank{}. （结果用数值表示）
\end{problem}

\begin{problem}
赌博有陷阱.某种赌博每局的规则是：赌客先在标记有1,2,3,4,5的卡片中随机摸取一张，将卡片上的数字作为其赌金（单位：元）；随后放回该卡片，再随机摸取两张，将这两张卡片上数字之差的绝对值1.4倍作为其奖金（单位：元）. 若随机变量 \(\xi_{1}\) 和 \(\xi_{2}\) 分别表示赌客在一局赌博中的赌金和奖金，则 \(E(\xi_{1}) - E(\xi_{2}) =\)\fillinblank{}（元）.
\end{problem}

\begin{problem}
已知函数 \( f(x) = \sin x \). 若存在 \(x_1\)，\(x_2\)，\(\cdots\)，\(x_m\) 满足 \( 0 \leqslant x_1 < x_2 < \cdots < x_m \leqslant 6\pi \). 且 \( |f(x_1) - f(x_2)| + |f(x_2) - f(x_3)| + \cdots + |f(x_{m-1}) - f(x_m)| = 12  \)（\(m \geqslant 2\)，\(m \in \mathbb{N}^*\)），则 \( m \) 的最小值为\fillinblank{}.
\end{problem}

\begin{problem}
在锐角三角形 \(ABC\) 中，\(\tan A = \frac{1}{2}\)，\(D\) 为边 \(BC\) 上的点，\(\triangle ABD\) 与 \(\triangle ACD\) 的面积分别为 2 和 4. 过 \(D\) 作 \(DE \perp AB\) 于 \(E\)，\(DF \perp AC\) 于 \(F\)，则 \(\overrightarrow{DE} \cdot \overrightarrow{DF} =\fillinblank{}\).
\end{problem}

\section{单选题}

\begin{problem}
设 \(z_{1}\)，\(z_{2} \in \mathbf{C}\)，则“\(z_{1}\)，\(z_{2}\) 中至少有一个数是虚数”是“\(z_{1} - z_{2}\) 是虚数”的（\quad）

\choices
    {充分非必要条件}
    {必要非充分条件}
    {充要条件}
    {既非充分又非必要条件}
\end{problem}

\begin{problem}
已知点 \(A\) 的坐标为 \((4\sqrt{3}, 1)\)，将 \(OA\) 绕坐标原点 \(O\) 逆时针旋转 \(\frac{\pi}{3}\) 至 \(OB\)，则点 \(B\) 的纵坐标为（\quad）

\choices
    {\(\frac{3\sqrt{3}}{2}\)}
    {\(\frac{5\sqrt{3}}{2}\)}
    {\(\frac{11}{2}\)}
    {\(\frac{13}{2}\)}
\end{problem}

\begin{problem}
记方程\circled{1}： \( x^{2} + a_{1}x + 1 = 0 \) ，方程\circled{2}： \( x^{2} + a_{2}x + 2 = 0 \) ，方程\circled{3}： \( x^{2} + a_{3}x + 4 = 0 \) ，其中 \(a_{1}\)，\(a_{2}\)，\(a_{3}\) 是正实数. 当 \(a_{1}\)，\(a_{2}\)，\(a_{3}\) 成等比数列时，下列选项中，能推出方程\circled{3}无实根的是（\quad）

\choices
    {方程\circled{1}有实根，且\circled{2}有实根}
    {方程\circled{1}有实根，且\circled{2}无实根}
    {方程\circled{1}无实根，且\circled{2}有实根}
    {方程\circled{1}无实根，且\circled{2}无实根}
\end{problem}

\begin{problem}
设 \(P_{n}(x_{n}, y_{n})\) 是直线 \(2x - y = \frac{1}{n + 1} (n \in \mathbf{N}^{*})\) 与圆 \(x^{2} + y^{2} = 2\) 在第一象限的交点，则极限 \(\lim\limits_{n \to \infty} \frac{y_{n} - 1}{x_{n} - 1} =\)（\quad）

\choices
    {\(-1\)}
    {\(-\frac{1}{2}\)}
    {1}
    {2}
\end{problem}

\section{解答题}

\begin{problem}
如图，在长方体 \(ABCD - A_{1}B_{1}C_{1}D_{1}\) 中，\(AA_{1} = 1\)，\(AB = AD = 2\)，\(E\)，\(F\) 分别是棱 \(AB\)，\(BC\) 的中点. 证明 \(A_{1}\)，\(C_{1}\)，\(F\)，\(E\) 四点共面，并求直线 \(CD_{1}\) 与平面 \(A_{1}C_{1}FE\) 所成角的正弦值.

\bitmapfigure[max width=0.62\linewidth,max totalheight=5.8cm]{img/2015/shanghai_science/q19_fig1.png}
\end{problem}

\begin{problem}
如图，\(A\)，\(B\)，\(C\) 三地有直道相通，\(AB = 5\) 千米，\(AC = 3\) 千米，\(BC = 4\) 千米. 现甲，乙两警员同时从 \(A\) 地出发匀速前往 \(B\) 地，过 \(t\) 小时，他们之间的距离为 \(f(t)\) （单位：千米）. 甲的路线是 \(AB\)，速度为 \(5\) 千米/小时，乙的路线是 \(ACB\)，速度为 \(8\) 千米/小时. 乙到达 \(B\) 地后在原地等待. 设 \(t = t_1\) 时，乙到达 \(C\) 地.
\begin{enumerate}
\item 求 \( t_{1} \) 与 \( f(t_{1}) \) 的值；
\item 已知警员的对讲机的有效通话距离是 3 千米. 当 \( t_{1} \leqslant t \leqslant 1 \) 时，求 \( f(t) \) 的表达式，并判断 \( f(t) \) 在 \( [t_{1}, 1] \) 上的最大值是否超过 \(3\). 说明理由.
\end{enumerate}

\bitmapfigure[max width=0.36\linewidth,max totalheight=4.8cm]{img/2015/shanghai_science/q20_fig1.png}
\end{problem}

\begin{problem}
已知椭圆 \(x^{2} + 2y^{2} = 1\) ，过原点的两条直线 \(l_{1}\) 和 \(l_{2}\) 分别与椭圆交于点 \(A\)，\(B\) 和 \(C\)，\(D\). 记得到的平行四边形ACBD的面积为S.
\begin{enumerate}
\item 设 \(A(x_{1}, y_{1})\)，\(C(x_{2}, y_{2})\). 用 \(A\)，\(C\) 的坐标表示点 \(C\) 到直线 \(l_{1}\) 的距离，并证明 \(S = 2|x_{1}y_{2} - x_{2}y_{1}|\).
\item 设 \(l_{1}\) 与 \(l_{2}\) 的斜率之积为 \(-\frac{1}{2}\)，求面积 \(S\) 的值.
\end{enumerate}
\end{problem}

\begin{problem}
已知数列 \(\{a_{n}\}\) 与 \(\{b_{n}\}\) 满足 \(a_{n + 1} - a_n = 2(b_{n + 1} - b_n)\)，\(n \in \mathbf{N}^*\).
\begin{enumerate}
\item 若 \(b_{n} = 3n + 5\)，且 \(a_{1} = 1\)，求 \(\{a_{n}\}\) 的通项公式；
\item 设 \(\{a_{n}\}\) 的第 \(n_0\) 项是最大项，即 \(a_{n_0} \geqslant a_n\) (\(n \in \mathbf{N}^*\)). 求证： \(\{b_n\}\) 的第 \(n_0\) 项是最大项；
\item 设 \(a_1 = \lambda < 0\)，\(b_n = \lambda^n (n \in \mathbf{N}^*)\). 求 \(\lambda\) 的取值范围，使得 \(\{a_n\}\) 有最大值 \(M\) 与最小值 \(m\)，且 \(\frac{M}{m} \in (-2, 2)\).
\end{enumerate}
\end{problem}

\begin{problem}
对于定义域为 \(\R\) 的函数 \(g(x)\). 若存在正常数 \(T\)，使 \(\cos g(x)\) 是以 \(T\) 为周期的函数，则称 \(g(x)\) 为余弦周期函数，且称 \(T\) 为其余周期. 已知 \(f(x)\) 是以 \(T\) 为余弦周期的余弦周期函数，其值域为 \(\R\). 设 \(f(x)\) 单调递增，\(f(0) = 0\)，\(f(T) = 4\pi\).
\begin{enumerate}
\item 验证 \( h(x) = x + \sin \frac{x}{3} \) 是以 \( 6\pi \) 为余弦周期的余弦周期函数；
\item 设 \( a < b \)，证明对任意 \( c \in [f(a), f(b)] \)，存在 \( x_0 \in [a, b] \)，使得 \( f(x_0) = c \)；
\item 证明：“\( u_0 \) 为方程 \( \cos f(x) = 1 \) 在 \([0, T]\) 上的解”的充要条件是“\( u_0 + T \) 为方程 \( \cos f(x) = 1 \) 在 \([T, 2T]\) 上的解”，并证明对任意 \( x \in [0, T] \) 都有 \( f(x + T) = f(x) + f(T) \).
\end{enumerate}
\end{problem}
